\section{A uniform bias bound for an equal bulk}
\label{sec:equal-bulk-bias}

The global atomic argument requires control of all possible biases after
conditioning on exceptional coordinates. The next theorem supplies
this control, including the second-order error needed to rule out a
single exceptional coordinate.

For $r\ge1$, define
\[
 B_r(a,b)=\E_X\left[\chi_{[r]}(X)
           \tanh\left(a\sum_{i=1}^rX_i+b\right)\right].
\]
The normalization $q_n(a)$ refers to the positive balanced coefficient
of odd degree $n$.

\begin{theorem}[An arbitrary bias cannot improve an equal bulk]
\label{thm:equal-bulk-bias}
Uniformly for $a\in[1,2]$, the following hold for all sufficiently
large $r$.
\begin{enumerate}
\item If $r$ is odd, then
\[
 \sup_{b\in\R}|B_r(a,b)|=q_r(a),
\]
and the unique maximizing bias is $b=0$.
\item If $r$ is even, then
\[
 \sup_{b\in\R}|B_r(a,b)|
 \le q_{r+1}(a)\,[1+O(r^{-2})].
\]
Every maximizing bias lies at $b=a+O(r^{-1})$ or
$b=-a+O(r^{-1})$.
\end{enumerate}
All constants and the sufficiently-large threshold are independent of
$a$.
\end{theorem}

\begin{proof}
We first prove that every maximizing bias lies in a fixed compact
interval. This step is needed before using a local expansion around
the periodic maxima.

Put
\[
 K_a(x)=\frac{\tanh(x+a)-\tanh(x-a)}2>0.
\]
The kernel and every fixed derivative decrease exponentially in
$|x|$, uniformly for $a\in[1,2]$. Taking one selected finite difference
first gives the exact identity
\begin{equation}\label{eq:bias-local-kernel}
 B_r(a,b)=\E_{X\in\{-1,1\}^{r-1}}
       \chi_{[r-1]}(X)K_a\left(a\sum_iX_i+b\right).
\end{equation}
Its alternating periodization is
\begin{equation}\label{eq:local-kernel-periodization}
 \sum_{j\in\mathbb Z}(-1)^jK_a(c+2aj)=T_a(c+a).
\end{equation}
Indeed $K_a$ is the difference of two shifted tanh functions, and
$T_a(c-a)=-T_a(c+a)$.

Write $N=r-1$. Let $p_j$ be the probabilities of
$\frac12\sum_{i=1}^NX_i$, on the integer lattice when $N$ is even
and on the half-integer lattice when $N$ is odd. Put
$p_*=\max_jp_j$. Up to a common sign,
\eqref{eq:bias-local-kernel} is
\[
 \sum_j(-1)^{j-j_0}p_jK_a(b+2aj),
\]
where $j_0$ is any fixed element of the mass lattice. Choose an index
$l$ of that lattice nearest $-b/(2a)$ and put $c=b+2al\in[-a,a]$.
After shifting the summation index, the coefficient is, up to a common
sign,
\begin{equation}\label{eq:shifted-binomial-kernel}
 \sum_{j\in\mathbb Z}(-1)^jp_{l+j}K_a(c+2aj).
\end{equation}

We record the binomial estimates used below. For a sufficiently small
absolute $\varepsilon>0$ and $|l|\le\varepsilon N$,
\begin{align}
 p_l/p_*&\le\exp\left[-c_1(|l|-1)_+^2/N\right],
 \label{eq:binomial-envelope}\\
 \sum_j\left|\frac{p_{l+j}}{p_l}-1\right|
       |K_a(c+2aj)|&\le C\frac{|l|+1}{N}.
 \label{eq:binomial-local-variation}
\end{align}
The constants apply to both lattice parities. To check the second
estimate directly, use the consecutive ratio
\[
 \frac{p_{s+1}}{p_s}=\frac{N/2-s}{N/2+s+1}.
\]
Summing logarithms for $|l|,|j|\le\varepsilon N$ gives
\[
 \left|\frac{p_{l+j}}{p_l}-1\right|
 \le C\frac{|l||j|+j^2+|j|}{N}
 \exp\left(C\frac{|l||j|+j^2+|j|}{N}\right).
\]
Choose $\varepsilon$ so that the last exponential is absorbed by
the fixed exponential decay of $K_a$. For the remaining indices,
$p_l\ge p_*e^{-C\varepsilon^2N}$ and the kernel tail is
$O(e^{-c\varepsilon N})$; decreasing $\varepsilon$ absorbs their
ratio and proves \eqref{eq:binomial-local-variation}.
The first estimate follows from the same consecutive ratios, using
$\log(1-x)\le-x$.

The periodic maximum $|T_a|\le F_0(a)$, together with
\eqref{eq:shifted-binomial-kernel}--\eqref{eq:binomial-local-variation},
now gives
\begin{equation}\label{eq:bias-envelope}
 \frac{|B_r(a,b)|}{p_*}
 \le e^{-c_1(|l|-1)_+^2/N}
       \left[F_0(a)+C\frac{|l|+1}{N}\right]
 \qquad(|l|\le\varepsilon N).
\end{equation}
For $|l|>\varepsilon N$, split
\eqref{eq:shifted-binomial-kernel} at $|j|=|l|/2$.
The first part uses the Gaussian binomial bound
$p_{l+j}\le p_*e^{-cl^2/N}$, and the second uses the exponentially
small kernel tail. The result is exponentially small in $N$;
the same argument applies if $l$ is outside the binomial support.

There is a central candidate of size
$p_*[F_0(a)-O(N^{-1})]$: use $b=0$ for odd $r$ and $b=a$ for even
$r$. Since $F_0$ is uniformly positive on $[1,2]$,
\eqref{eq:bias-envelope} confines every maximizing index to
$|l|\le L$ for an absolute constant $L$. More explicitly, for
$L\le|l|\le\sqrt N$, the negative quadratic term $-cl^2/N$
dominates the linear correction $C(|l|+1)/N$ once $L$ is large.
For $\sqrt N\le|l|\le\varepsilon N$, reduce $\varepsilon$ so
that the fixed Gaussian loss dominates that correction. All remaining
indices were already exponentially small. Thus every maximizing bias
lies in a fixed compact interval, uniformly in $a$ and $r$.

On that compact interval, expand the binomial probabilities in
\eqref{eq:bias-local-kernel} around the center. The localized kernel
allows the polynomial-in-$j$ remainders to be summed with any fixed
number of bias derivatives. This gives a uniform expansion through
$O(r^{-2})$, after central-binomial normalization. Its leading response
is $T_a(b+a)$ for odd $r$ and $T_a(b)$ for even $r$, up to a common
sign. Their maxima are nondegenerate. Consequently every maximizing
bias lies within $O(r^{-1})$ of one of finitely many periodic peaks,
and optimizing within such a neighborhood changes its normalized
value by only $O(r^{-2})$.

For odd $r$, the peaks are $b=2la$. The Fourier identity for
$\cos(2l\theta)$ and the balanced ratio
\[
 \frac{q_{r+2j}(a)}{q_r(a)}=1-\frac jr+O(r^{-2})
\]
give, uniformly over the fixed range of $l$,
\begin{equation}\label{eq:odd-bias-peak-loss}
 \frac{|B_r(a,2la)|}{q_r(a)}
       =1-\frac{2l^2}{r}+O(r^{-2}).
\end{equation}
One way to see the coefficient is to normalize the polynomial
$\cos(2l\theta)$ so that its value at
$y=\sin^2\theta=1$ is one; its derivative there is $2l^2$.
Every nonzero $l$ loses even after the $O(r^{-2})$ local improvement.
At $l=0$, the response is even in $b$, so zero is exactly stationary;
strict local concavity makes it the unique local, hence global,
maximum of the absolute coefficient.

For even $r$, set $n=r+1$. The peaks are $b=(2l+1)a$.
The identity $B_r(a,b)=C_n(a,\ldots,a,b)$ identifies them with
one odd harmonic in an odd-degree parity atom. The sine polynomial
therefore gives
\begin{equation}\label{eq:even-bias-peak-loss}
 \frac{|B_r(a,(2l+1)a)|}{q_n(a)}
       =1-\frac{2l(l+1)}n+O(n^{-2}).
\end{equation}
Only $l=0,-1$, corresponding to $b=\pm a$, survive.
Local optimization shifts them by $O(r^{-1})$ and increases their
relative value by $O(r^{-2})$, as required.
\end{proof}
